\section*{Capítulo \ref{ch:g11:vect} --- Vetores e retas no plano}

\begin{solution}{exo:g11:vect:1}
$\vect{AB}\,(3, 4)$, $\vect{AC}\,(-3, 2)$,
$\vect{AB} + \vect{AC}\,(0, 6)$,
$2\vect{AB} - 3\vect{AC}\,(6 + 9,\ 8 - 6) = (15, 2)$. O \omterm{prop:g10:coordgeom:midpoint}{ponto médio} de
$[BC]$ é $\left(\frac{5 + (-1)}{2}, \frac{3+1}{2}\right) = (2, 2)$.
\end{solution}

\begin{solution}{exo:g11:vect:2}
$\det = 4 \times 3 - (-2) \times (-6) = 12 - 12 = 0$: \omterm{def:g11:vect:collinear}{colineares} (de
fato, $\vec v = -\frac12 \vec u$).

$\det = 2 \times 7 - 3 \times 5 = -1 \neq 0$: não são \omterm{def:g11:vect:collinear}{colineares}.

$\det = 0 \times (-8) - 0 \times 3 = 0$: \omterm{def:g11:vect:collinear}{colineares} (ambos verticais).
\end{solution}

\begin{solution}{exo:g11:vect:3}
Por $A(2,1)$, direção $(3, -1)$: identificar $(-b, a) = (3, -1)$ dá
$a = -1$, $b = -3$, logo $-x - 3y + c = 0$; substituindo $A$:
$-2 - 3 + c = 0$, $c = 5$. Multiplicando por $-1$: $x + 3y - 5 = 0$.

Por $B(0,4)$ e $C(2,0)$: $\vect{BC}\,(2, -4)$, logo $a = -4$, $b = -2$:
$-4x - 2y + c = 0$; substituindo $B$: $-8 + c = 0$, $c = 8$.
Simplificando por $-2$: $2x + y - 4 = 0$. Confira com $C$:
$4 + 0 - 4 = 0$.
\end{solution}

\begin{solution}{exo:g11:vect:4}
Um \omterm{def:g11:vect:direction}{vetor diretor} é $(-b, a) = (2, 3)$. \omterm{def:g10:numbers:interunion}{Interseções}: $x = 0$ dá $y = 3$,
ponto $(0, 3)$; $y = 0$ dá $x = -2$, ponto $(-2, 0)$. Para $P(2,6)$:
$3 \times 2 - 2 \times 6 + 6 = 0$, logo $P \in d$. Para $Q(1,4)$:
$3 - 8 + 6 = 1 \neq 0$, logo $Q \notin d$.
\end{solution}

\begin{solution}{exo:g11:vect:5}
$\vect{AB}\,(3, 6)$ e $\vect{AC}\,(-3, -6)$:
$\det = 3 \times (-6) - (-3) \times 6 = 0$, logo $A$, $B$, $C$ são
\omterm{def:g11:vect:collinear}{colineares}.

$\vect{DE}\,(2, 3)$ e $\vect{DF}\,(5, 8)$:
$\det = 2 \times 8 - 5 \times 3 = 1 \neq 0$: $D$, $E$, $F$ não são
\omterm{def:g11:vect:collinear}{colineares}.
\end{solution}

\begin{solution}{exo:g11:vect:6}
Primeiro par: somando as duas \omterm{def:g10:algebra:equation}{equações},
$(2x + y - 7) + (x - y + 1) = 3x - 6 = 0$, logo $x = 2$ e então
$y = x + 1 = 3$: elas se encontram em $(2, 3)$ (verificação:
$4 + 3 - 7 = 0$).

Segundo par: $ab' - a'b = 1 \times 4 - (-2) \times (-2) = 0$, logo as
retas são paralelas; multiplicar $e$ por $-2$ dá $-2x + 4y - 6 = 0$, que
difere de $e'$ ($-6 \neq 1$): estritamente paralelas, sem \omterm{def:g10:numbers:interunion}{interseção}.
\end{solution}

\begin{solution}{exo:g11:vect:7}
Para $d: 2x - y + 3 = 0$: direção $(-b, a) = (1, 2)$, e o ponto $(0, 3)$
está em $d$, logo $d = \{(t,\ 3 + 2t) : t \in \R\}$.

Para a reta paramétrica que passa por $(1, -3)$ com direção $(2, 1)$:
$a = 1$, $b = -2$, logo $x - 2y + c = 0$; substituindo $(1, -3)$:
$1 + 6 + c = 0$, $c = -7$. \omterm{def:g10:algebra:equation}{Equação}: $x - 2y - 7 = 0$ (confira com
$t = 1$, ponto $(3, -2)$: $3 + 4 - 7 = 0$).
\end{solution}

\begin{solution}{exo:g11:vect:8}
Uma reta paralela tem o mesmo par $(a, b)$: $4x - 3y + c = 0$; passando
por $P(1, -2)$: $4 + 6 + c = 0$, logo $c = -10$ e a reta é
$4x - 3y - 10 = 0$.

O paralelismo de $mx + 2y - 5 = 0$ com $d$ exige
$ab' - a'b = 4 \times 2 - m \times (-3) = 8 + 3m = 0$, logo
$m = -\frac83$.
\end{solution}

\begin{solution}{exo:g11:vect:9}
O \omterm{prop:g10:coordgeom:midpoint}{ponto médio} de $[AB]$ é
$M\!\left(\frac{-1+3}{2}, \frac{0+2}{2}\right) = (1, 1)$. A \omterm{def:g10:stats:median}{mediana} é a
reta $(CM)$ com $C(1, -4)$: os dois pontos têm \omterm{def:g10:coordgeom:system}{abscissa} $1$, de modo que
a \omterm{def:g10:stats:median}{mediana} é a reta vertical
\[
x - 1 = 0 .
\]
\end{solution}

\begin{solution}{exo:g11:vect:10}
$\det(\vec u, \vec v) = m(m - 1) - 3 \times 2 = m^2 - m - 6$. Ele se
anula quando $m^2 - m - 6 = 0$: $\Delta = 1 + 24 = 25$, raízes
$m = \frac{1 \pm 5}{2}$, isto é, $m = 3$ ou $m = -2$.
\end{solution}

\begin{solution}{exo:g11:vect:11}
\emph{1.} $I = \left(\frac12, 0\right)$. Então
$\vect{DI} = \left(\frac12 - 0,\ 0 - 1\right) = \left(\frac12,
-1\right)$ e $\vect{DJ} = \frac23\vect{DI} = \left(\frac13,
-\frac23\right)$, logo
\[
J = D + \vect{DJ} = \left(\frac13,\ \frac13\right).
\]

\emph{2.} $\vect{AJ}\left(\frac13, \frac13\right)$ e
$\vect{AC}\,(1, 1)$:
$\det = \frac13 \times 1 - 1 \times \frac13 = 0$, logo $A$, $J$, $C$ são
\omterm{def:g11:vect:collinear}{colineares} --- na verdade, $\vect{AJ} = \frac13\vect{AC}$: a reta $(DI)$
corta a diagonal $(AC)$ exatamente a um terço do seu comprimento a
partir de $A$.
\end{solution}

\begin{solution}{pb:g11:vect:1}
\textbf{1.} $x = 1 + 3t$, $y = 2 - t$. Da segunda,
$t = 2 - y$; substituindo: $x = 1 + 3(2 - y) = 7 - 3y$:
\omterm{def:g10:algebra:equation}{equação} cartesiana $x + 3y - 7 = 0$.

\textbf{2.} Para $ax + by + c = 0$, um \omterm{def:g11:vect:direction}{vetor diretor} é
$(-b, a)$: aqui $(5, 2)$. Um ponto: $y = 1$ dá $x = 1$:
$(1, 1)$. Paramétrica: $x = 1 + 5t$, $y = 1 + 2t$.

\textbf{3.} De $x = 7 - 3y$ em $2x - 5y + 3 = 0$:
$14 - 6y - 5y + 3 = 0$, logo $y = \frac{17}{11}$ e
$x = 7 - \frac{51}{11} = \frac{26}{11}$: o ponto
$\left(\frac{26}{11}, \frac{17}{11}\right)$.

\textbf{4.} Direções $(5, 2)$ e $(3, -1)$:
$\det = 5 \times (-1) - 2 \times 3 = -11 \neq 0$: concorrentes,
como se encontrou. Para o segundo par: as direções $(5, 2)$ e
$(10, 4)$ têm $\det = 0$ (paralelas), e dobrar
$2x - 5y + 3 = 0$ dá $-4x + 10y - 6 = 0 \neq -4x + 10y + 1 = 0$:
estritamente paralelas.

\textbf{5.} $\vect{AB}(3, 2)$, $\vect{AC}(9, 6)$:
$\det = 18 - 18 = 0$: \omterm{def:g11:vect:collinear}{colineares}.

\textbf{6.} $P$: de $x = 2t$ e $y = 3 + t$ vem
$y = 3 + \frac x2$. $Q$: de $x = 10 - t$ vem $t = 10 - x$, logo
$y = 2(10 - x) - 1 = 19 - 2x$. Cruzamento:
$3 + \frac x2 = 19 - 2x$ dá $x = 6.4$, $y = 6.2$: as
trajetórias se cruzam em $(6.4,\ 6.2)$.

\textbf{7.} $P$ passa ali quando $2t = 6.4$: $t = 3.2$ h. $Q$,
quando $10 - t = 6.4$: $t = 3.6$ h. Vinte e quatro minutos de
diferença: nenhuma colisão. Aviso: uma carta mostra
\emph{trajetórias}; só os relógios paramétricos dizem se dois
movimentos ocupam o mesmo ponto no mesmo instante.

\textbf{8.} $D^2(t) = (3t - 10)^2 + (4 - t)^2 =
10t^2 - 68t + 116$: vértice em $t = \frac{68}{20} = 3.4$ h,
onde $D^2 = 0.4$: \omterm{def:g10:numbers:approx}{aproximação} mínima
$D = \sqrt{0.4} \approx 0.63$ milha em $t = 3.4$.

\textbf{9.} Violada: $0.63 < 1$. $D^2(t) < 1$ lê-se
$10t^2 - 68t + 115 < 0$: $\Delta = 4624 - 4600 = 24$, raízes
$\frac{68 \pm \sqrt{24}}{20} = 3.4 \pm 0.245$: os navios ficam
próximos demais de $t \approx 3.16$ a $t \approx 3.64$ ---
cerca de $29$ minutos de infração. Um deles precisa mudar de
rota.

\textbf{10.} Cada coordenada de cada navio é uma \omterm{def:g11:func:function}{função}
\emph{afim} de $t$, de modo que cada diferença de \omterm{def:g10:coordgeom:system}{coordenadas} é
afim, e $D^2$ --- uma soma de dois quadrados de \omterm{def:g10:reffunc:affine}{funções afins}
--- é do segundo grau (com coeficiente dominante não negativo).
A \omterm{def:g10:numbers:approx}{aproximação} mínima, portanto, é sempre lida em um vértice.

\textbf{11.} Encontrar $P$ em $t = 3$, isto é, em $(6, 6)$:
$(3a, 3b) = (6, 6)$ dá $(a, b) = (2, 2)$, velocidade
$\sqrt{8} \approx 2.8$ nós.

\textbf{12.} $D(0, 1)$, $I\left(\frac12, 0\right)$: direção
$\left(\frac12, -1\right)$, isto é, $(1, -2)$: \omterm{def:g10:algebra:equation}{equação}
$2x + y = 1$.

\textbf{13.} $(AC)$: $y = x$. Então $2x + x = 1$:
$x = \frac13$: o ponto $\left(\frac13, \frac13\right)$, a um
terço do caminho de $A$ até $C$.

\textbf{14.} $B(1, 0)$, $K\left(\frac12, 1\right)$: direção
$\left(-\frac12, 1\right)$, \omterm{def:g10:algebra:equation}{equação} $2x + y = 2$. Com $y = x$:
$x = \frac23$: o ponto $\left(\frac23, \frac23\right)$. As duas
cevianas cortam $(AC)$ em seus dois pontos de trissecção:
teorema demonstrado, duas vezes três linhas de aritmética.

\textbf{15.} Todo paralelogramo é levado ao quadrado unitário
escolhendo-se $A$ como origem e $\vect{AB}$, $\vect{AD}$ como
unidades dos dois eixos. Essa escolha preserva tudo aquilo de
que o teorema fala --- \omterm{prop:g10:coordgeom:midpoint}{pontos médios}, colinearidade e razões ao
longo de uma mesma reta (tudo exprimível com \omterm{def:g11:vect:vector}{vetores} e
escalares) ---, de modo que demonstrar o enunciado no quadrado o
demonstra para todo paralelogramo. (Ela \emph{não} transportaria
comprimentos nem ângulos: esses o referencial distorce.)

\textbf{16.} $E\left(1, \frac12\right)$: reta $(AE)$:
$y = \frac x2$. Diagonal $(DB)$: de $(0,1)$ a $(1,0)$:
$y = 1 - x$. \omterm{def:g10:numbers:interunion}{Interseção}: $\frac x2 = 1 - x$: $x = \frac23$: o
ponto $\left(\frac23, \frac13\right)$ --- de novo um ponto de
trissecção, agora da outra diagonal. As cevianas dos \omterm{prop:g10:coordgeom:midpoint}{pontos
médios} adoram terços.

\textbf{17.} As duas primeiras: $x + y = 4$ e $2x - y = 2$ dão
$3x = 6$: $(2, 2)$. A terceira: $2 - 4 = -2$: satisfeita.
Concorrentes em $(2, 2)$.

\textbf{18.} $mx - y + 2 - 3m = m(x - 3) + (2 - y) = 0$: para a
\omterm{def:g10:algebra:equation}{equação} valer para \emph{todo} $m$, tome $x = 3$, $y = 2$ --- e,
de fato, toda $d_m$ contém $(3, 2)$. Um feixe de retas por um
único ponto, uma reta para cada \omterm{def:g10:reffunc:affine}{coeficiente angular} $m$.

\textbf{19.} Paralela a $y = 2x + 7$: \omterm{def:g10:reffunc:affine}{coeficiente angular}
$m = 2$, a reta $2x - y - 4 = 0$. Passando pela origem:
$2 - 3m = 0$: $m = \frac23$.

\textbf{20.} Cartesiana: substitua um ponto e a pertinência
está respondida. Reduzida: \omterm{def:g10:reffunc:affine}{coeficiente angular} e \omterm{def:g10:reffunc:affine}{coeficiente
linear} à vista, \omterm{def:g10:functions:graph}{gráfico} imediato. Paramétrica: o relógio já
embutido --- colisões, \omterm{def:g10:numbers:approx}{aproximação} mínima, interceptação. O
\omterm{def:g11:vect:det}{determinante}: uma multiplicação em cruz e o paralelismo está
decidido. O referencial escolhido: um paralelogramo vira um
quadrado unitário, e um teorema clássico de trissecção vira três
\omterm{def:g10:numbers:interunion}{interseções} de retas.
\end{solution}
